\chapter{INTRODUCTION}\label{sec:intro}
\section{THIS IS A FIRST LEVEL HEADING}
Recent advances in nanoscale transport have been summarized by  \cite{Smith2021Chap}, who provided a detailed theoretical framework for phonon drag effects.

For further background on non-linear oscillations, see the discussion by Jones \cite{Jones2018Sec}.

Experimental validation of metamaterial cloaking at microwave frequencies was reported by \cite{Lee2019Conf}.

Monte Carlo studies of lattice gauge models at finite temperature were explored by \cite{Garcia2022Master}.

Nguyen’s doctoral work \cite{Nguyen2020PhD} established a correlation between lattice strain and superconducting transition temperature.

The velocity field data used for model verification were retrieved from the Dryad repository \cite{Hansen2017Data}.

Neural-network modeling was implemented using the open-source TensorFlow library \cite{TensorFlow2015}.

A concise overview of the fundamental postulates can be found on Wikipedia \cite{Wikipedia2025Quantum}.

Environmental compliance in laboratory procedures followed ISO 14001:2015 \cite{ISO14001:2015}.

A broad philosophical overview of foundational issues is collected in an Oxford handbook \cite{Batterman2013OxfordHandbook}.

For a focused chapter inside a book, see Penrose’s discussion of spinors \cite{Penrose1987Inbook}.

Deep residual networks were introduced at \cite{He2016ResNetConference}.

The full ICML 2020 volume is available as an edited proceedings from PMLR \cite{ICML2020Proceedings}.

Shannon’s master’s thesis established the algebraic basis of switching circuits \cite{Shannon1937Masters}.

Turing’s Princeton dissertation laid groundwork for ordinal logics \cite{Turing1938PhD}.

Strain data for GW150914 can be obtained from LOSC \cite{LOSCGW150914Dataset}.

Many numerical examples here were produced with NumPy \cite{NumPySoftware}.

Physical constants are retrieved from NIST’s online database \cite{NISTCODATAOnline}.

Time formatting follows the ISO date/time standard \cite{ISO8601Standard}.

\subsection{This is a second level heading}
Classical mechanics describes the motion of a particle of mass \(m\)
under a force \(\mathbf{F}\) through Newton’s second law:
\begin{equation}
  \mathbf{F} = m \mathbf{a}.
  \label{eq:newton}
\end{equation}
Equation~\eqref{eq:newton} remains a cornerstone of physics \cite{Newton1687Principia}.

For waves on a string, the standard one-dimensional wave equation is
\begin{equation}
  \frac{\partial^2 y}{\partial t^2} = v^2 \frac{\partial^2 y}{\partial x^2},
  \label{eq:wave}
\end{equation}
where \(v = \sqrt{T/\mu}\) is the wave speed determined by tension \(T\)
and linear density \(\mu\) \cite{French1971Vibrations}.

In quantum mechanics, the time-independent Schrödinger equation reads
\begin{equation}
  -\frac{\hbar^2}{2m}\nabla^2 \psi + V\psi = E\psi,
  \label{eq:schrodinger}
\end{equation}
first introduced by Schrödinger in 1926 \cite{Schrodinger1926Quantisation}.

\section{FIGURE EXAMPLES}

Scientific documents often include figures to illustrate results, experimental setups,
or conceptual frameworks. Figure~\ref{fig:single} shows a single image with a caption.
The convention for captions and references follows guidelines in \cite{KopkaDalyGuide}.

\begin{figure}[h!]
  \centering
  \includegraphics[width=0.75\linewidth]{example-image-a}
  \caption{A single placeholder example image.}
  \label{fig:single}
\end{figure}

Multiple related panels can be combined as subfigures (Fig.~\ref{fig:subfigures});
see general practices for composite figures in \cite{Tufte2001VisualDisplay}.

\begin{figure}[h]
  \centering
  \begin{subfigure}{0.45\linewidth}
    \includegraphics[width=\linewidth]{example-image-b}
    \caption{First subfigure.}
  \end{subfigure}\hfill
  \begin{subfigure}{0.45\linewidth}
    \includegraphics[width=\linewidth]{example-image-c}
    \caption{Second subfigure.}
  \end{subfigure}
  \caption{Example of two related subfigures side by side.}
  \label{fig:subfigures}
\end{figure}

Larger works often combine multiple figures in a sequence to show
progressive results (Fig.~\ref{fig:sequence}), following presentation techniques
from \cite{EdwardRHeathLaTeXIllustrations}.

\begin{figure}[h]
  \centering
  \includegraphics[width=0.3\linewidth]{example-image-a}
  \includegraphics[width=0.3\linewidth]{example-image-b}
  \includegraphics[width=0.3\linewidth]{example-image-c}
  \caption{A sequence of related figures illustrating progressive results.}
  \label{fig:sequence}
\end{figure}

Finally, figures may also appear in floating environments with wide layouts
such as a two-column article (not shown here), or may include multiple panels
referenced collectively (e.g., Figs.~\ref{fig:single}–\ref{fig:sequence}).

\section{TABLE EXAMPLES}

Experimental data often appear in tabular form.
Table~\ref{tab:constants} lists selected fundamental constants
recommended by CODATA \cite{MohrCODATA2018}. These values are widely
used in numerical simulations and physical calculations.

\begin{table}[h!]
\centering
\caption{Selected fundamental physical constants (CODATA 2018) \cite{MohrCODATA2018}.}
\label{tab:constants}
\begin{tabular}{lcc}
\toprule
\textbf{Quantity} & \textbf{Symbol} & \textbf{Value} \\
\midrule
Speed of light & $c$ & $2.9979\times10^{8}\ \text{m/s}$ \\
Planck constant & $h$ & $6.6261\times10^{-34}\ \text{J·s}$ \\
Elementary charge & $e$ & $1.6022\times10^{-19}\ \text{C}$ \\
Boltzmann constant & $k_{\mathrm{B}}$ & $1.3806\times10^{-23}\ \text{J/K}$ \\
\bottomrule
\end{tabular}
\end{table}

\bigskip

Table~\ref{tab:materials} shows sample material properties adapted from
standard references on solid-state materials \cite{Callister2018Materials}.
Such data are frequently used to compare thermal and mechanical performance.

\begin{table}[h!]\footnotesize
\centering
\caption{Typical properties of selected engineering materials \cite{Callister2018Materials}.}
\label{tab:materials}
\begin{tabular}{lccc}
\toprule
\textbf{Material} & \textbf{Density (g/cm$^3$)} & \textbf{Elastic Modulus (GPa)} & \textbf{Thermal Conductivity (W/m·K)} \\
\midrule
Aluminum & 2.70 & 69 & 237 \\
Copper   & 8.96 & 110 & 401 \\
Steel (AISI 1020) & 7.87 & 200 & 51 \\
Silicon  & 2.33 & 130 & 149 \\
\bottomrule
\end{tabular}
\end{table}

\bigskip

In social-science or management contexts, smaller textual tables are also common.
Table~\ref{tab:survey} illustrates a Likert-scale example as used in
survey analysis \cite{Likert1932Technique}.

\begin{table}[h!]
\centering
\caption{Example of a Likert-scale question and responses \cite{Likert1932Technique}.}
\label{tab:survey}
\begin{tabular}{lccccc}
\toprule
\textbf{Statement} & \textbf{SD} & \textbf{D} & \textbf{N} & \textbf{A} & \textbf{SA} \\
\midrule
Physics is enjoyable & 3 & 8 & 15 & 20 & 4 \\
Mathematics is easy  & 5 & 10 & 12 & 18 & 5 \\
\bottomrule
\end{tabular}
\end{table}

Table~\ref{tab:survey} uses the abbreviations:
SD (Strongly Disagree), D (Disagree), N (Neutral),
A (Agree), SA (Strongly Agree).

\section{SECTIONING AND CROSS-REFERENCING}\label{sec:sectioning}
This section demonstrates how to create hierarchical headings and cross-references.  
Always assign each major part of your thesis a logical structure:
\begin{itemize}
  \item \texttt{\textbackslash section\{...\}}
  \item \texttt{\textbackslash subsection\{...\}}
  \item \texttt{\textbackslash subsubsection\{...\}}
\end{itemize}

For example, the methodology details are explained in Section~\ref{sec:methods}, 
while the results are summarized in Section~\ref{sec:results}.

\subsection{Methodology Example}
\label{sec:methods}

The study followed a two-stage process: simulation and verification.  
The simulation parameters are listed in Table~1 (not shown here).  
Cross-referencing is done using commands like \texttt{\textbackslash ref\{sec:results\}}.

\subsection{Results Example}
\label{sec:results}

This subsection provides the findings, which are discussed further in Section~\ref{sec:discussion}.

\subsubsection{Statistical Analysis}
Descriptive and inferential statistics were used to interpret data.

\section{DISCUSSION}
\label{sec:discussion}

Section~\ref{sec:sectioning} described the use of sections and labels; 
here we connect those references logically in the narrative.

% ==================================================
\section{LISTS}
\label{sec:lists}

Lists help organize information clearly.  
LaTeX supports both numbered and bulleted lists, as shown below.

\subsection{Numbered List (Enumerate)}
Use numbered lists for procedural steps:
\begin{enumerate}
  \item Review existing literature.
  \item Define the problem statement.
  \item Design and conduct experiments.
  \item Analyze data and interpret results.
  \item Draw conclusions and propose future work.
\end{enumerate}

\subsection{Bulleted List (Itemize)}
Use bulleted lists for unordered information:
\begin{itemize}
  \item Experimental setup
  \item Numerical model
  \item Analytical solution
\end{itemize}

You can also nest lists if necessary:
\begin{itemize}
  \item Input parameters:
    \begin{enumerate}
      \item Temperature
      \item Pressure
      \item Volume
    \end{enumerate}
  \item Output:
    \begin{enumerate}
      \item Energy
      \item Entropy
    \end{enumerate}
\end{itemize}

% ==================================================
\section{INLINE MATHEMATICS}
\label{sec:inline-math}

Mathematics can appear either inline or displayed separately.  
Inline math is enclosed within dollar signs \texttt{\$...\$}, 
and it flows naturally with the text.

For example, the relationship between distance $d$, velocity $v$, 
and time $t$ is expressed as $d = v t$.  
Similarly, the kinetic energy is given by $E_k = \tfrac{1}{2} m v^2$, 
and the potential energy by $E_p = m g h$.  

For larger expressions, use display mode:
\[
E = \sqrt{(p c)^2 + (m_0 c^2)^2}
\]
Equation numbering can also be added with the \texttt{equation} environment.

% ==================================================
\section{QUOTATIONS AND FOOTNOTES}
\label{sec:quotes-footnotes}

Proper quotation formatting improves readability and avoids plagiarism.  
Short quotations are included inline, for example:  
According to Einstein, \enquote{Imagination is more important than knowledge.}

For longer excerpts, use the \texttt{quote} environment:
\begin{quote}
Science is a way of thinking much more than it is a body of knowledge.  
\hfill---Carl Sagan
\end{quote}

You may also add explanatory or bibliographic comments using footnotes.  
For instance, quantum mechanics was formulated in the early twentieth century,\footnote{Historically, the development involved both matrix mechanics (Heisenberg, 1925) and wave mechanics (Schrödinger, 1926).}  
which revolutionized physics.

Footnotes should be concise and supplementary, never containing major arguments or references that belong in the bibliography.

